\begin{trapbox}

	\begin{description}[style=nextline, leftmargin=2.8em, labelsep=0.5em, itemsep=0.35em]
		\item[\textbf{\textcolor{CTrap}{易错点一（排列组合混淆）}}]
			求解“从3个元素中选2个排成一列”时，错误地使用组合公式 $C_3^2 = 3$。
		\item[\textbf{\textcolor{CTrap}{正确理解（区分有序无序）：}}]
			排列与顺序有关，正确应为 $A_3^2 = 3\times 2 = 6$。
		\item[\textbf{\textcolor{CTrap}{错因分析（顺序忽视）：}}]
			未考虑元素顺序，误把排列问题当成组合问题。
	\end{description}

	\medskip
	\begin{description}[style=nextline, leftmargin=2.8em, labelsep=0.5em, itemsep=0.35em]
		\item[\textbf{\textcolor{CTrap}{易错点二（公式记忆错误）}}]
			将排列数公式错写为 $A_n^m = \dfrac{n!}{m!(n-m)!}$。
		\item[\textbf{\textcolor{CTrap}{正确理解（分母仅阶乘）：}}]
			排列数公式分母只有 $(n-m)!$，没有 $m!$，即 $A_n^m = \dfrac{n!}{(n-m)!}$。
		\item[\textbf{\textcolor{CTrap}{错因分析（组合公式干扰）：}}]
			与组合数公式 $C_n^m = \dfrac{n!}{m!(n-m)!}$ 混淆所致。
	\end{description}

	\medskip
	\begin{description}[style=nextline, leftmargin=2.8em, labelsep=0.5em, itemsep=0.35em]
		\item[\textbf{\textcolor{CTrap}{易错点三（忽略定义域）}}]
			计算 $A_5^6$ 时直接套用公式，或认为 $A_5^0$ 无意义。
		\item[\textbf{\textcolor{CTrap}{正确理解（条件确保）：}}]
			必须满足 $n,m \in \mathbb{N}$ 且 $m \le n$；$m=0$ 时 $A_n^0 = 1$。
		\item[\textbf{\textcolor{CTrap}{错因分析（前提缺失）：}}]
			未检查 $m \le n$ 或对 $0!$ 定义不清楚。
	\end{description}

\end{trapbox}
